\section{Plant Phenotyping}

\paragraph{}
Abiotic and biotic stressors limit cultivar biomass and yield \cite{grimeEvidenceExistenceThree1997}, whilst human population growth strains the global food supply. Plant breeding, the process of inducing desirable, inheritable traits in plants through controlled human action \cite{ortonIntroduction2020}, is now crucial in addressing the challenge of providing an adequate global food supply. Plant breeders require effective techniques to produce germplasm for cultivars that maximise yield and biomass under various adverse environmental conditions. In the last few decades, image processing techniques in plant phenotyping have proven to be an effective tool for plant breeding \cite{anandBreedingTechniquesDispense2023}.

\subsection{What Is Plant Phenotyping?}

\paragraph{}
Plant phenotyping is the evaluation of phenotypes in plants \cite{fioraniFutureScenariosPlant2013}, where phenotypes are primitive or complex traits of an organism that can be observed or quantified and evaluated numerically \cite{johannsenGenotypeConceptionHeredity1911a}.  Phenotypes can alternatively be seen as a combination of all morphological, physiological, chemical, and behavioural aspects that represent the individual plant \cite{guoPhenotypingPlants2014, mohammedReviewPlantPhenotyping2025, liReviewImagingTechniques2014}. 

\subsection{Plant Phenotyping As Method For Improving Plant Traits}

\paragraph{}
Phenotypical traits serve as proxies in determining whether specific genotypes have favourable traits in particular environmental conditions. Quantitative data produced by plant phenotyping is most helpful in analyzing plant-environment interactions in plant breeding \cite{fioraniFutureScenariosPlant2013} for the development of germplasm with improved plant traits \cite{liReviewImagingTechniques2014}. The data produced by phenotyping is used as a proxy for assessing a plant's performance in its environment. For example, stressors or disturbances in plants are induced by environmental conditions to which organisms are exposed, and these stressors induce symptoms that can be detected and evaluated using phenotyping. Hence, phenotypical traits of a plant have a complex dynamic feedback loop with their environment \cite{walterPlantPhenotypingBean2015}. Many phenotypical traits are observable or quantifiable by image processing techniques.

\subsection{Image-Obtainable Phenotypical Traits As Indicators Of Plant Health, Growth, And Stress}

\subsubsection{Simple Taxonomy of Phenotypical Traits}
\paragraph{}
A phenotypical trait is classifiable as morphological, physiological, or temporal. Additionally, the literature indicates morphological and physiological phenotypical traits are classifiable as holistic-based or component-based \cite{daschoudhuryLeveragingImageAnalysis2019}. Holistic properties consider the plant as a whole, whereas component-based features examine individual parts of the organism, such as its fruit, flower, leaves, stem, etc. 

\subsubsection{Leaf-Derived Phenotypical Traits: Robust Proxy Indicators of Plant Metabolic and Physiological Status}
\paragraph{}
Leaf-derived traits are among the most used indicators of plant status because leaves rapidly respond to environmental stimuli and robustly indicate various factors that determine plant health, growth, and stress \cite{zhangHighthroughputPhenotypingPlant2023}. Traits of leaves can monitor and predict a host of plant health, growth and stress factors. For example, leaf area predicts plant productivity because reduced leaf area results in decreased chlorophyll content and less photosynthesis. Therefore, leaf area is a good proxy measure for evaluating plant growth and has been used to monitor and estimate crop yield and health \cite{zhangLeafAreaIndex2021}. Various other leaf morphological traits exist in plant phenotyping. These include leaf number and inclination angle for monitoring plant growth and development \cite{praveenkumarRosettePlantSegmentation2020, zhangHighthroughputPhenotypingPlant2023}, and leaf thickness for measuring abiotic stressors like water scarcity \cite{afzalLeafThicknessPredict2017}. Additionally, physiological traits of leaves exist that indicate various factors of plant status, such as stomatal conductance for detecting biotic stressors like pathogens \cite{ishimweApplicationsThermalImaging2014} and leaf water content for drought stress \cite{watanabeFundamentalStudyWater2022}. Considering the above, leaf-derived phenotypical traits are robust proxies of various metabolic and physiological plant parameters.

\subsubsection{Alternative Phenotypical Traits}
\paragraph{}
Although leaves play an important role in evaluating plant growth, health and resilience to stress, these evaluations cannot be scaled to the whole plant \cite{zhangHighthroughputPhenotypingPlant2023}. Rather, other parts of the plant should be phenotyped to understand the processes governing plant status comprehensively. The literature is rife with such examples that study traits of roots \cite{bodnerHyperspectralImagingNovel2018}, stems \cite{daschoudhuryAutomatedStemAngle2017}, seeds \cite{zhangHighThroughputPhenotypingSeed2018}, and fruit \cite{abebeImageBasedHighThroughputPhenotyping2023}.